\chapter{Einführung}
\label{ch:Einfuehrung}

Passantenfrequenzen weisen wiederkehrende zeitliche und räumliche Muster auf.
Typischerweise unterscheiden sich Tag- und Nachtstunden, Werktage und
Wochenenden sowie stark und schwach frequentierte Standorte. Abweichungen vom
zeit- und standortabhängig erwarteten Verhalten werden in dieser Arbeit als
Anomaliekandidaten betrachtet. Sie sind nicht automatisch mit einem fachlich
bestätigten realen Ereignis gleichzusetzen. Vollständig fehlende Übertragungen
werden getrennt als Datenqualitätsereignisse behandelt.

Datengrundlage ist ein API-Snapshot des Smart-City-Portals der Stadt Bad
Nauheim mit Messungen von 30~Sensorstandorten im halboffenen Zeitraum vom
30.~Juni 2025, 00:00~Uhr, bis zum 1.~Juli 2026, 00:00~Uhr UTC
\cite{badnauheim_pedestrian_2026}. Der globale Zeitraum umfasst 366 Tage; die
einzelnen Sensorreihen besitzen jedoch unterschiedliche Messbeginne und
teilweise umfangreiche Lücken.

\section{Motivation}
\label{sec:Einfuehrung_Motivation}

Eine automatisierte Erkennung auffälliger Verläufe kann die manuelle Prüfung
umfangreicher Sensordaten unterstützen. Erschwert wird sie durch fehlende
vollständige reale Anomaliekennzeichnungen und ein Normalverhalten, das mit
Tageszeit, Wochentag und Standort variiert. Ein einheitlicher statischer
Grenzwert kann diese Unterschiede nur eingeschränkt berücksichtigen
\cite{hundman2018spacecraft}.

\section{Zielsetzung}
\label{sec:Einfuehrung_Zielsetzung}

Ziel der Arbeit ist der Vergleich von vier getrennten Verfahren zur Erkennung
auffälliger Passantenfrequenzen. Quantil-LSTM, Quantil-Transformer und
LSTM-Autoencoder werden ohne reale Anomalieklassen selbstüberwacht trainiert;
eine robuste saisonale $t-24$-Methode dient als Baseline. Die Alarmregeln
werden im Kalibrierungsabschnitt mit kontrollierten synthetischen Anomalien
ausgewählt und anschließend mit identischen Kriterien verglichen.

Die zentrale Forschungsfrage lautet: Wie unterscheiden sich die vier Verfahren
hinsichtlich der punkt- und ereignisbasierten Erkennung der fünf definierten
synthetischen Anomaliearten in den stündlichen Passantenfrequenzen von Bad
Nauheim?

Bewertet wird damit die technische Erkennungsfähigkeit für definierte
Veränderungen; belastbare Erkennungsraten für unbekannte reale Anomalien lassen
sich ohne vollständige reale Labels nicht ableiten.

Die folgenden Kapitel behandeln Grundlagen, Stand der Technik, Methode,
Evaluation und Ergebnisse; eine Zusammenfassung schließt die Arbeit ab.
